\reportsection{14}{Acquisition}{Acquisition: Fetching Upstream, Reworked for the Official .deb}

\unithead{Mechanism}

\reppar{\code{scripts/setup/detect-host.sh}'s \code{detect\_architecture} is a \code{case} on \code{uname -m} (\code{x86\_64}/\code{aarch64}) that sets \code{architecture} (\code{amd64}/\code{arm64}) and, in the same branch, hardcodes a versioned \code{downloads.claude.ai/releases/win32/\{x64,arm64\}/<ver>/Claude-<sha>.exe} URL, \code{claude\_exe\_sha256}, and \code{claude\_exe\_filename}. There is no runtime resolution in the build itself: the URL and hash are literal strings, refreshed only by CI. Anything outside \code{x86\_64}/\code{aarch64} hard-exits. \code{scripts/setup/download.sh}'s \code{download\_claude\_installer} either copies a local file (the \code{--exe} escape hatch) or \code{wget}s the hardcoded URL, then calls \code{verify\_sha256()} and exits on mismatch. The \code{.exe} is \code{7z x}-tracted twice, once for the NSIS wrapper and once for the nested \code{.nupkg}, with \code{version} parsed from the nupkg filename via:}

\begin{codeblock}
grep -oP 'AnthropicClaude-\K[0-9]+\.[0-9]+\.[0-9]+(?=-full|-arm64-full)'
\end{codeblock}

\reppar{A daily cron in \code{.github/workflows/check-claude-version.yml} runs \code{scripts/resolve-download-url.py} (Playwright, headless Chromium) against the \code{claude.ai} redirect endpoints, since the bare redirect sits behind a Cloudflare bot challenge a plain \code{curl}/\code{wget} cannot pass. The resolved URL and version are diffed against \code{detect-host.sh}; on a change the workflow \code{sed}-rewrites the URL/hash lines, recomputes SHA-256 (hex) and Nix SRI (base64) for both architectures, updates \code{nix/claude-desktop.nix}'s \code{version}/\code{hash} fields, commits, sets \code{CLAUDE\_DESKTOP\_VERSION}, and tags and pushes \code{v\{REPO\_VERSION\}+claude\{CLAUDE\_VERSION\}}, the push that triggers the release build. \code{nix/claude-desktop.nix} itself is a \code{fetchurl} per \code{system} (\code{x86\_64-linux}/\code{aarch64-linux}) against the same URLs, each pinned with a Nix SRI hash: one resolver, two consumers.}

\unithead{Origin}

\reppar{The true origin is \code{d8f4bbe} (2024-12-26, the project's initial commit): \code{build-deb.sh} hardcoded a single unversioned \code{storage.googleapis.com} URL, \code{wget}-fetched with no SHA-256 verification at all and no version automation. Every later piece, arch-aware URLs, hash verification, CI-driven bumping, Nix pins, was added afterward, not present at inception. The redirect approach had a same-day false start on 2025-11-28: \code{5c5eb39} (\#143) switched to the \code{claude.ai/api/desktop/.../redirect} endpoint with a spoofed-browser \code{curl} header set to dodge Cloudflare, and the very next commit, \code{1405e1c}, reverted both changes the same day, back to the static URL and plain \code{wget}. The static-URL approach eventually rotted for real: \#151's \code{3e813c5} (2026-01-05) replaced it with Playwright-based resolution after the URLs were found "stuck at v0.14.10," the exact Cloudflare problem \#143 had failed to solve with header spoofing five weeks earlier. SHA-256 verification did not arrive until \#310's \code{1fc4e83} (2026-03-20), which added \code{verify\_sha256()} and wired it into both the installer and Node tarball fetches; before that commit a corrupted or substituted download would silently proceed into the build.}

\begin{chronology}
2024-12 & n/a & \code{d8f4bbe}: origin; hardcoded, unversioned URL in \code{build-deb.sh}; no hash check, no automation. \\
2025-08 & closes \#104 & \code{a987938}: \code{check-claude-version.yml} created, daily cron, \code{CLAUDE\_DESKTOP\_VERSION} var, tag scheme. \\
2025-11 & \#143 & \code{5c5eb39}: switched to the redirect endpoint plus spoofed-header \code{curl}; \code{1405e1c} reverted both the same day. \\
2026-01 & \#151; \#146 & \code{3e813c5}: Playwright-based \code{resolve-download-url.py} added, fixing URLs stuck at v0.14.10; same day, \#146 added the \code{--exe} escape hatch. \\
2026-02/03 & \#266 & \code{ff9fd3d} \arrow{} \code{55f7b27}: \code{nix/claude-desktop.nix} created with per-arch \code{fetchurl} plus SRI pins. \\
2026-03 & \#310 & \code{1fc4e83}: \code{verify\_sha256()} added; installer and Node tarball verified after download. \\
2026-04 & \#443 & \code{ff4821e} plus \code{564f465}: \code{build.sh} split into \code{scripts/setup/detect-host.sh} plus \code{download.sh}; CI paths updated in the same PR. \\
2026-04 & \#535 & \code{4cc63bf}: workflow third-party actions pinned to commit SHAs, post trivy-action supply-chain incident. \\
2026-07 & n/a & \code{9f3820c}: Rebase Phase 0; \code{official-deb.sh} written, pinned pool paths plus SHA-256, not yet wired into \code{build.sh}. \\
2026-07 & n/a & \code{d9cef9e}: Rebase Phase 1; \code{build.sh} sources \code{official-deb.sh}; legacy download chain deleted, Nix reduced to a stub. \\
2026-07 & n/a & \code{cafb4cc}: Rebase Phase 5; \code{check-claude-version.yml} rewritten against the official Packages index (no Playwright); \code{mirror-official-deb} job added. \\
\end{chronology}

\methodbanner{\textbf{Fate under the rebase.} \brework{} The acquisition target changed from a Windows \code{.exe} installer to the official Linux \code{.deb}, forcing every piece of the chain to change with it, though the shape (arch-detect, pinned URL/hash, download, verify, extract, daily CI bump, Nix pin) stayed the same. \code{scripts/setup/official-deb.sh} replaces \code{download.sh}, \code{resolve-download-url.py}, and \code{fetch-electron-binary.js} outright (commit \code{d9cef9e}: "Deleted: download.sh, resolve-download-url.py, fetch-electron-binary.js, setup\_electron\_asar, scripts/staging/*"). The official \code{.deb} is fetched from \code{OFFICIAL\_APT\_BASE=downloads.claude.ai/claude-desktop/apt/stable}, a plain-HTTPS APT pool with no bot challenge, so Playwright is gone: \code{resolve\_official\_deb()} is a \code{curl} plus \code{awk} parse of the Packages index. It remains SHA-256 pinned by two independent mechanisms: build-time constants checked through the same \code{verify\_sha256()} helper \#310 added, and CI-time reads straight from the Packages index. Extraction is deliberately \code{ar}/\code{tar} rather than \code{dpkg-deb}, so rpm-family and Arch hosts can build too. A new \code{mirror-official-deb} CI job uploads the pinned \code{.deb} for both architectures to our own GitHub Releases as insurance against the upstream pool rotating a version out. \code{nix/claude-desktop.nix} is an honest stub for now: it \code{throw}s toward the old derivation in git history, with @typedrat named as the pending owner, and the version-bump workflow already detects the stub and skips the SRI step rather than failing. \\[3pt]\textbf{Affects.} Arch-specific rather than distro-specific (amd64/arm64 download URLs); the Nix SRI hashes are the one format-specific piece. DE-agnostic.}
